\chapter{Introduction}
\label{ch:introduction}

\section{Background}
A \textbf{network} (or graph) is a set of nodes connected by links, and is one of the most general
ways of representing relationships in data. In many real systems, however, the same entities are
related in several different ways simultaneously: two researchers may collaborate in physics and,
separately, in mathematics. Such systems are modelled as \textbf{multiplex (multilayer)
networks}~\cite{kivela2014}, in which every type of relationship forms its own \emph{layer} over a
shared set of nodes, so that the same node keeps the same identity in every layer. Link prediction
in such networks is an active and rapidly evolving research area~\cite{zangari2024,heuristickdd2024}, with
growing use of \textbf{attention-based graph neural networks}~\cite{gatv2,attnsurvey2023}.

\textbf{Link prediction} estimates, for a pair of nodes that are not currently connected, how likely
they are to become connected. It underlies collaboration and friend recommendation, the completion
of partially observed biological and knowledge networks, and the analysis of how scientific fields
grow.

\section{Motivation}
Not all links are equally easy or equally valuable to predict. Pairs that already share many
neighbours are easy. The most valuable links, however, are frequently the \textbf{weak-tie links} --
pairs drawn from different communities that share no common neighbours yet. Granovetter's theory of
the strength of weak ties~\cite{granovetter1973} argues that precisely these bridging connections
carry unique, non-redundant information. This creates a difficulty: classical methods score a pair
from the number of neighbours it shares, so on a weak-tie link that number is zero, the score is
zero, and the prediction is effectively random. A second limitation is that most methods analyse
each layer in isolation and discard the information that a node's behaviour in one layer carries for
another. These two gaps motivate this work.

\section{Objectives}
The internship pursued three objectives: (i) to study link prediction in multiplex networks and
review the literature, with emphasis on \emph{attention mechanisms} in graph neural networks; (ii)
to implement and fairly evaluate a representative set of baseline methods; and (iii) to design an
\textbf{attention-based model} that predicts weak-tie links, combines information across all layers,
and remains interpretable.

\section{Contributions}
The central contribution is an attention-based graph neural network that unifies two attention
mechanisms -- a per-layer, weak-tie-aware graph-attention encoder and a cross-layer attention that
weighs the layers -- together with a fusion gate. The model attains the highest ROC-AUC and a
statistically significant improvement over the neural baselines, predicts weak-tie links where
classical methods are random, and exposes interpretable layer-importance weights. This report is
extended into the master's thesis, where each part is presented in greater depth.
